%! TeX root: ../article.tex
% Resumen y palabras clave en español. Se \input dentro del entorno abstract
% de article.tex (llncs ya imprime el encabezado "Abstract", por eso aquí no se
% incluye ningún título ni configuración de página).

El seguimiento manual de grafiti no autorizado se vuelve inviable a medida que
crecen los bancos de imágenes. Este trabajo diseña, implementa y evalúa
experimentalmente un sistema modular para agrupar por similitud visual un
dataset no etiquetado de fotografías de grafiti, con el objetivo de
caracterizar tanto la calidad del agrupamiento como el coste computacional de
sus componentes en función del tamaño del banco. La arquitectura encadena
etapas intercambiables (segmentación, extracción de características,
almacenamiento, reducción de dimensionalidad y agrupamiento) y un marco de
experimentación recorre combinaciones de las mismas midiendo, por ejecución,
tiempos por etapa, memoria y métricas de calidad. La configuración base
resultante (cabeza estilística de DINOv2 con \textit{fine-tuning} + UMAP-10 +
HDBSCAN + imagen completa) se valida tanto sobre el dataset de $6416$ imágenes
de grafitis como sobre subconjuntos disjuntos etiquetados manualmente. La
introducción de una cabeza clasificadora para separar estilos artísticos separa
estos con un índice de Rand ajustado de 0.722 (K-medias con $k{=}4$ conocido
\textit{a priori}), mientras que un agrupamiento por autoría se queda cerca del
azar, evidenciando la insuficiencia de datos anotados por autor en el conjunto
de prueba. El estudio de escalabilidad muestra que la búsqueda por similitud en
almacenamiento con SQLite y la ingesta son los cuellos de botella dominantes a
esta escala y que otros factores como la complejidad de los algoritmos de
reducción o agrupamiento empezarían a serlo a escalas bastante mayores. Se
concluye que las métricas intrínsecas de evaluación del agrupamiento pueden
inducir a error y que la elección del extractor de características debe
anclarse en evidencia supervisada.

\keywordname{} Grafiti $\cdot$ Agrupamiento de imágenes $\cdot$
\textit{Embeddings} profundos $\cdot$ Reducción de dimensionalidad $\cdot$
HDBSCAN $\cdot$ Búsqueda por similitud $\cdot$ Coste computacional
